\section{Guarded native probes and the admissible-intermediate counter-bridge}
\label{sec:136-admissible-native-probes}

\paragraph{Archive, inputs, and scope.}
This concerns cote136, physical pp.~117 and119, footers116
and118, scans117 and119, and the common-source and homotopy
questions on physical p.~121, footer120, scan121.  We use the
checked internal statements ARCH-136-DERIVED-TAIL-COMPARISON,
ARCH-136-RELATIVE-COHERENT-MOORE,
ARCH-136-TAIL-EXT-VANISHING, and the actual free presentation of
ARCH-136-RELATIVE-CELL-PATHS.  The specified finite-path homotopy
is ARCH-136-HOM-DEFINITION.  The nonzero example is
ARCH-136-CRITICAL-NATIVE-PATH-COMPONENT.  No additional outside
theorem or preprint is used.  This is a counter-bridge for that
particular route to a homotopy obstruction, not a decision of
the archive's homotopy question over \(\mathbb Z\).

Work over \(S=\Gamma_{\mathbb Z}(T)\).  Fix an odd
\(q\geq1\), an integer \(p\geq2\), and a positive native
weight \(m\).  Put
\[
               D=p(q+1),\qquad v=pm,\qquad r=D-2.
                                                               \tag{136.NP1}
\]
Choose a free graded resolution \(P_\bullet\to I_M\), with
all generator weights at least \(M\).  For a simplicial pair
\(Y\subseteq X\), let
\[
 E_M(X,Y)=S\amalg_{\mathcal N_M(Y)}\mathcal N_M(X),
                                                               \tag{136.NP2}
\]
where the boundary map to \(S\) sets every native generator
and every coherent parameter to zero.  Thus the boundary here
is the augmentation boundary.  It is not an arbitrary previously
prescribed nonzero boundary under an old algebra.

For a target \(C\), abbreviate
\[
 \mathcal H_C(X,Y)=\operatorname{Tot}
   \operatorname{Hom}^{0}_{S,\mathrm{gr}}
                    (P_\bullet[v],C(X,Y)),\qquad
 \mathcal P_C(X,Y)=P_p(\mathcal H_C(X,Y)).
                                                               \tag{136.NP3}
\]
The differentials are the full-resolution differentials of
(136.DT3) and (136.DT15).  A strict map \(E_M(X,Y)\to C\)
is exactly a native section \(a\in C^{m,q}(X,Y)\) and a
parameter cochain \(\beta=(-V,-U)\in\mathcal P_C(X,Y)^r\)
satisfying
\[
 d\beta=\omega_p(a),\qquad
 \omega_p(a)=\bigl(f_{a\gamma_{p-1}(da)},f_{\gamma_p(da)}\bigr).
                                                               \tag{136.NP4}
\]
The maps \(f\) multiply the zeroth free-resolution component
by its augmentation and vanish in later components.  The next
resolution-index equations in (136.NP4) are precisely both
literal bottom relations of the strict PD presentation.

\paragraph{ARCH-136-GUARDED-NATIVE-PROBE.}
Let \(C\) be admissible, assume its native row
\(C_*^{m,q}\) is horizontally acyclic, and require
\[
                         M+v\geq W_C(D).                \tag{136.NP5}
\]
Then finite-path homotopy classes of strict maps (136.NP2) to
\(C\), with their augmentation boundary retained, have a
canonical bijection
\[
 [E_M(X,Y),C]_s\ \cong\ H^r(X,Y;S^v/pS^v).
                                                               \tag{136.NP6}
\]
It sends a zero-native map with closed parameter cochain \(\beta\)
to its canonical cone cohomology class.  Every map is homotopic
to such a zero-native map.  The bijection is natural in admissible
targets for which the stated guards hold, and in maps of pairs
and augmentation-compatible maps of coefficient-tail resolutions.
Primality is not needed.

\begin{proof}
The derived-tail comparison and the coefficient-tail Ext theorem
give the canonical isomorphism
\[
              H^r\mathcal P_C(X,Y)
                       \cong H^r(X,Y;S^v/pS^v).          \tag{136.NP7}
\]
Only the adjacent degrees \(r,r+1\) are needed for the cone
comparison.  The stronger bound (136.NP5) also covers every
extension and relative native primitive below.  All comparisons
are induced by the coefficient map (136.DT7); hence they commute
with strict coefficient-preserving theory maps, as well as with
the pair and free-resolution maps stated in that theorem.

First reduce a given map \((a,\beta)\) to zero native input.
Put \(Z=X\times\Delta[1]\) and
\[
        Y_0=(Y\times\Delta[1])\cup(X\times\{0\}).
\]
The native row guard extends zero at endpoint zero and \(a\)
at endpoint one to a section \(\widetilde a\) on \(Z\)
which vanishes on \(Y_0\).  This is the checked relative
free-chain row-extension argument.  The pair \((Z,Y_0)\)
is contractible: the interval contraction to zero preserves its
subspace and ends inside that subspace.  Therefore the relative
derived-tail comparison and the cone sequence make
\(H^{D-1}\mathcal P_C(Z,Y_0)=0\).  Choose
\[
        \widetilde\beta\in\mathcal P_C(Z,Y_0)^r,
                      \qquad d\widetilde\beta
                                        =\omega_p(\widetilde a).
                                                               \tag{136.NP8}
\]
This is a full coherent parameter lift, zero at the initial
endpoint.  Write \(\beta_*\) for its other endpoint and set
\(\alpha=\beta-\beta_*\).  It is a closed cochain in
\(\mathcal P_C(X,Y)^r\).  The native interpolation
\(\widetilde a\), with parameters
\(\widetilde\beta+\operatorname{pr}_X^*\alpha\), is an
actual strict path from \((0,\alpha)\) to \((a,\beta)\).
The linear parameter equations and both bottom relations hold;
the free strict PD universal property supplies all products and
divided powers.  No integral contraction of a nonflat parameter
module has been assumed.

Next consider two zero-native maps \((0,\alpha_0)\) and
\((0,\alpha_1)\).  If their cone classes agree, write
\(\alpha_1-\alpha_0=d\gamma\).  Extend the cochain
with endpoint values zero and \(\gamma\) to a cochain
\(\widetilde\gamma\) on the interval cylinder, relative
to \(Y\times\Delta[1]\).  This extension is coordinatewise:
the resolution modules are free and every relevant form row has
weight at least \(M+v\) in the acyclic range.  Arbitrary
products of these surjections are surjective.  The zero native
input with closed parameters
\[
       \operatorname{pr}_X^*\alpha_0+d\widetilde\gamma
                                                               \tag{136.NP9}
\]
gives a strict one-edge path with the desired endpoints.

Conversely suppose a finite oriented interval path \(J\)
connects those two zero-native maps.  Its native input vanishes
on
\[
        Y'=(Y\times J)\cup(X\times\partial J)
                               \ \subseteq\ Z'=X\times J.
\]
ARCH-136-RELATIVE-COHERENT-MOORE supplies a relative cochain
\(\rho\in\mathcal P_C(Z',Y')^r\) with differential equal
to this path's native pair.  Subtract \(\rho\) from its
actual parameter cochain.  The result is closed relative to
\(Y\times J\) and still has endpoints \(\alpha_0,\alpha_1\).
Projection to \((X,Y)\) is a homotopy equivalence of pairs.
The derived-tail and cone comparisons therefore identify the
two endpoint maps on cohomology: both are inverses of that
projection's isomorphism.  Thus
\([\alpha_0]=[\alpha_1]\).  This argument covers every
integer \(p\geq2\), not just primes.

Every closed cone class defines a zero-native strict map by
(136.NP4).  Reduction to zero native input proves surjectivity
of the classification, and the preceding necessity and sufficiency
prove injectivity and independence of every reduction choice.
Postcomposing a chosen reduction path gives a reduction in the
next target, and (136.NP7) is natural.  This proves all stated
naturality assertions.
\end{proof}

\paragraph{ARCH-136-ADMISSIBLE-INTERMEDIATE-COUNTERBRIDGE.}
Let \(A,C\) be admissible, let \(f_0,f_1:A\rightrightarrows C\)
be strict arrows, and let \(\theta:E_M(X,Y)\to A\) be a
strict map.  Assume only the original native-row guard in \(A\).
After restricting the coherent parameters to a sufficiently high
coefficient tail, the two composites \(f_0\theta,f_1\theta\)
are finite-path homotopic.  More precisely, choose \(L\geq M\)
with
\[
                L+v\geq\max\{W_A(D),W_C(D)\},          \tag{136.NP10}
\]
and a chain map from a free resolution of \(I_L\) to the
chosen resolution of \(I_M\), lifting the inclusion.  It
induces an actual strict map
\[
                  E_L(X,Y)\longrightarrow E_M(X,Y),
                                                               \tag{136.NP11}
\]
the identity on the native input.  The composites through
(136.NP11) are homotopic using at most three edges.  No native-row
guard in \(C\) is required for this positive conclusion.

In particular the nonzero prescribed pair of
ARCH-136-CRITICAL-NATIVE-PATH-COMPONENT cannot factor through
parallel admissible arrows with an intermediate whose native
row at \((m,q)\) is acyclic.  A countable family of those
pairs with fixed \(p,q\) and \(m_i\to\infty\) has no
admissible intermediate retaining all the prescribed pairs.
This is a statement about retaining specified maps under a
source, not the nonexistence of an admissible common source
for their targets.

\begin{proof}
A weight-preserving resolution comparison lifting
\(I_L\hookrightarrow I_M\) is obtained successively from
freeness and exactness.  Apply its components to every \(U_i,V_i\),
and leave \(a,b\) fixed.  It respects the native equations
because it commutes with the augmentation.  It respects the
bottom relations because it is a chain map also in the next
resolution index.  This proves the actual strict map (136.NP11).
Under the multiplier identification
\(\operatorname{Hom}_{S,\mathrm{gr}}(I_M[v],S)=S^v\),
restriction to \(I_L\) is the identity on \(S^v\).
The Ext vanishing and cone comparison show that it is likewise
the identity on the coefficient groups in (136.NP7).

Use the first part of the preceding proof in \(A\), after
(136.NP11), to connect \(\theta\) to a zero-native map
with closed parameter cochain \(\alpha\).  Postcomposition
gives paths from each \(f_j\theta\) to the zero-native map
with parameters \(f_j\alpha\).  The two latter cone classes
agree: the canonical comparison is induced by coefficients,
and both \(f_j\) preserve those coefficients identically.
The coordinatewise construction (136.NP9), using only parameter
row bounds in \(C\), supplies a zero-native one-edge path
between them.  Concatenating with the two postcomposed paths
gives the claimed path of at most three edges.  This explains
why no separate native-row hypothesis on \(C\) is needed.

In the critical sphere example, the endpoint parameter classes
are zero and \(T^{[v]}\bmod p\), and they remain different
under (136.NP11).  The prescribed maps are not joined by any
finite path, by that example's proved target guard and its
zero-native-loop argument.  The preceding conclusion is therefore
incompatible with the asserted admissible factorization.

For the countable family take the strict PD coproduct of the
individual zero-boundary native sources over \(S\), and
use their specified maps to the same standard target.  Every
expression has finite support, so these are actual strict maps.
For any alleged admissible intermediate, the fixed form degree
\(q\) has a finite native-row weight cutoff.  Some \(m_i\)
exceeds it.  Its parameter tail can be restricted further by
(136.NP11), without changing the nonzero class.  The just-proved
contradiction applies to that one factor.  For fixed \(p,q\),
one can equivalently choose an index which also exceeds the
finite parameter comparison bounds, since \(pm_i\to\infty\).
The targets in this example are the same standard theory, so
they certainly have a common admissible source.  The conclusion
excludes precisely the prescribed-intermediate construction.
\end{proof}

\subsection{Finite native degrees and the exact remaining homotopy issue}

There is also an explicit admissible denominator which removes
all exceptional low native weights in any fixed finite set of
positive form degrees.  For an admissible \(A\) and integer
\(K>0\), define \(A\langle K\rangle\) by keeping \(A\)
in weights at least \(K\), and keeping only \(S^w[0]\)
in smaller weights.  One precise construction, without assuming
injectivity of a low-weight coefficient map, is the strict
pullback
\[
 A\langle K\rangle
     =A\mathbin{\times}_{A/A_{\geq K}}(S/I_K).
                                                               \tag{136.NP12}
\]
The weight tails are differential PD ideals, so this is an
actual strict differential PD algebra.  In each low weight the
pullback is the coefficient copy; in each high weight it is
exactly \(A\).  All sufficiently high vertical comparisons
and horizontal rows are unchanged.  Thus
\(A\langle K\rangle\) is admissible and its projection to
\(A\) is a strict arrow.  This is the same weight-cutting
argument as the inherited AQ-136-WEIGHT-ONE construction; no
new model-structure assertion is involved.

For finitely many chosen positive form degrees \(q\), choose
\(K\) above their native-row cutoffs in \(A\).  The corresponding
rows of \(A\langle K\rangle\) are then acyclic in every
native weight: they are zero below \(K\), and are the guarded
old rows above it.  The counter-bridge applies to every such
zero-boundary probe after its own sufficiently high parameter-tail
restriction.  Consequently these probes cannot obstruct homotopy
after this one admissible precomposition in those fixed form
degrees.

This does not assemble the separate probe homotopies into a
homotopy on an arbitrary admissible source.  Full-boundary
attachments under an old algebra impose compatibility among
the cells; (136.NP2) fixes only their augmentation boundary.
Nor does a single finite \(K\) dominate arbitrary native-row
cutoffs in all form degrees.  The general homotopy reading (b)
over \(\mathbb Z\), and integral connectedness, remain open.

\paragraph{Verification status.}
The author completed a whole saved-proof reading, checking the
half-prism normalization, zero-native finite-path necessity and
sufficiency, every full-resolution equation, coefficient naturality,
the three-edge construction, tail restriction, and the strict
weight-cut pullback.  The classification and admissible-intermediate
argument were communicated to the parent and reviewer123 before
this module was written; those communications are not whole-file
reviews.  Reviewers123 and1348 separately communicated the
unbounded-family corollary before reviewing the saved proof;
that contribution is distinguished from any later whole-file check.
Root and reviewer123 have completed whole saved-proof readings with
PASS, including the weight-cut pullback and the exact nonassembly scope.
All new comparisons above use the stated uniform source and
parameter guards, not a silent assumption that the prescribed
nonadmissible example has an admissible completion.
